\lecture{10}{Knowledge Graph Embeddings}{Thursday, October 22, 2026}\label{ch:lec10}

% Local notation (lecture 10)
\providecommand{\vt}{\mathbf{t}}
\providecommand{\vg}{\mathbf{g}}
\providecommand{\C}{\mathbb{C}}
\providecommand{\cRe}{\operatorname{Re}}
\providecommand{\cIm}{\operatorname{Im}}
\providecommand{\Ker}{\operatorname{Ker}}
\providecommand{\cmark}{\textcolor{ForestGreen}{\checkmark}}
\providecommand{\xmark}{\textcolor{BrickRed}{$\boldsymbol{\times}$}}

% Local TikZ styles (prefix lec10)
\tikzset{
  lec10ent/.style={circle, fill=blue!80, minimum size=3.2mm, inner sep=0pt},
  lec10gen/.style={circle, fill=ForestGreen!85!black, minimum size=3.2mm, inner sep=0pt},
  lec10lbl/.style={font=\scriptsize\sffamily, fill=white, inner sep=1pt},
  lec10name/.style={font=\scriptsize\sffamily},
  lec10inf/.style={thick, RoyalBlue, -{Stealth[length=1.8mm]}},
  lec10gr/.style={thick, ForestGreen, -{Stealth[length=1.8mm]}},
  lec10pt/.style={circle, draw=black, fill=YellowOrange, minimum size=1.3mm, inner sep=0pt},
  lec10vec/.style={thick, -{Stealth[length=1.8mm]}},
  lec10ax/.style={thick, black!80},
  lec10proj/.style={densely dashed, black!60},
  lec10cap/.style={font=\scriptsize\sffamily, align=center},
}

\textbf{Outline.} Recap (heterogeneous graphs, RGCN) $\to$ knowledge graphs
(KGs) and KG datasets $\to$ KG completion with shallow embeddings $\to$ relation
patterns $\to$ TransE $\to$ TransR $\to$ DistMult $\to$ ComplEx $\to$ comparison
of the models and KG embeddings in practice.

% ---------------------------------------------------------------------
\section{Knowledge graphs}\label{sec:lec10-kg}

\subsection{Recap: heterogeneous graphs and relational GCN}\label{sec:lec10-recap}
\slides{3--4}

(Details in \cref{ch:lec09}.)
\begin{itemize}
  \item \textbf{Heterogeneous graphs:} a graph with \textbf{multiple relation
    types} (example input graph: nodes $A$--$F$, relations $r_1,r_2,r_3$, target
    node $A$).
  \item \textbf{Relational GCN:} learn from a graph with multiple relation types;
    \textbf{use different neural network weights for different relation types!}
    (computation graph of $A$: one neural network per relation type, then
    aggregation).
\end{itemize}

\subsection{Knowledge graphs and examples}\label{sec:lec10-kgdef}
\slides{5--9}

\begin{itemize}
  \item \term{Knowledge graph (KG)}: knowledge in graph form -- capture
    entities, types, and relationships:
    \begin{itemize}
      \item nodes are \textbf{entities}; nodes are labeled with their
        \textbf{types}; edges between two nodes capture \textbf{relationships}
        between entities (figure: entities $E_1,E_2,E_3$, each with $A_1,A_2$;
        relations $R_1,R_2,R_3$).
      \item \textbf{KG is an example of a heterogeneous graph.}
    \end{itemize}
  \item \textbf{Example: bibliographic networks.} Node types: paper, title,
    author, conference, year. Relation types: pubWhere, pubYear, hasTitle,
    hasAuthor, cite. Schema: Paper $\to$ Conference (pubWhere), Title
    (hasTitle), Author (hasAuthor), Year (pubYear), Paper (cite, self-loop).
  \item \textbf{Example: bio knowledge graphs.} Node types: drug, disease, adverse
    event, protein, pathways. Relation types: has\_func, causes, assoc, treats,
    is\_a. E.g.\ fulvestrant (drug) \textsc{treats} Breast Neoplasms (disease)
    and \textsc{causes} migraine (adverse event); Breast Neoplasms \textsc{is\_a}
    Neoplasms by Site; ESR1 (protein) \textsc{has\_func} response to estradiol
    (pathway); cerebral haemorrhage \textsc{assoc} CASP8 (protein).
  \item \textbf{Knowledge graphs in practice:} Google Knowledge Graph, Amazon
    Product Graph, Meta Graph API, IBM Watson, Microsoft Satori, Project
    Hanover/Literome, LinkedIn Knowledge Graph, Yandex Object Answer.
  \item \textbf{Applications -- serving information:} e.g.\ the Bing query
    ``latest films by the director of titanic'' returns ``Movies featuring James
    Cameron''.
\end{itemize}

\subsection{Knowledge graph datasets}\label{sec:lec10-datasets}
\slides{10--11}

\begin{itemize}
  \item \textbf{Publicly available KGs:} FreeBase, Wikidata, Dbpedia, YAGO, NELL,
    etc.
  \item \textbf{Common characteristics:} \textbf{Massive} -- millions of nodes and
    edges; \textbf{Incomplete} -- many true edges are missing.
  \item Given a massive KG, \textbf{enumerating all the possible facts is
    intractable!} $\Rightarrow$ \textbf{Can we predict plausible BUT missing
    links?}
  \item \textbf{Example: Freebase} -- $\sim$80 million \textbf{entities},
    $\sim$38K \textbf{relation types}, $\sim$3 billion \textbf{facts/triples}.
    Incomplete: \textbf{93.8\%} of persons from Freebase have no place of birth
    and \textbf{78.5\%} have no nationality! (Paulheim, 2017; Min et al., NAACL
    2013.)
  \item \textbf{Datasets: FB15k/FB15k-237} -- a \textbf{complete} subset of
    Freebase, used by researchers to learn KG models (\cref{tab:lec10-fb15k}).
\end{itemize}

\begin{table}[htbp]
  \centering\small
  \caption{FB15k and FB15k-237 (slide~11).}\label{tab:lec10-fb15k}
  \begin{tabular}{@{}lrrr@{}}
    \toprule
    Dataset & Entities & Relations & Total Edges\\
    \midrule
    FB15k     & 14,951 & 1,345 & 592,213\\
    FB15k-237 & 14,505 &   237 & 310,079\\
    \bottomrule
  \end{tabular}
\end{table}

% ---------------------------------------------------------------------
\section{Knowledge graph completion}\label{sec:lec10-completion}

\subsection{KG completion task and KG representation}\label{sec:lec10-task}
\slides{13--15}

\begin{itemize}
  \item \textbf{Given an enormous KG, can we complete the KG?} For a given
    \textbf{(head, relation)}, we predict missing \textbf{tails} (note: this is
    slightly different from the link prediction task).
  \item \textbf{Example task:} predict the tail ``Science Fiction'' for
    (``J.K. Rowling'', ``genre''), \cref{fig:lec10-completion}.
  \item \textbf{Recap: ``shallow'' encoding} (\cref{ch:lec02}): simplest encoding
    approach -- \textbf{encoder is just an embedding-lookup}; embedding matrix
    $\mZ$ has one column per node (the embedding vector for a specific node);
    number of rows = dimension/size of embeddings.
  \item \textbf{KG representation:} edges in KG are represented as
    \term{triples} $(h,r,t)$: \textbf{head} ($h$) has \textbf{relation} ($r$) with
    \textbf{tail} ($t$).
\end{itemize}

\begin{keyidea}
  Model entities and relations in embedding space $\R^k$: associate entities and
  relations with \textbf{shallow embeddings} (\textbf{we do not learn a GNN
  here!}). Given a triple $(h,r,t)$, the goal is that \textbf{the embedding of
  $(h,r)$ should be close to the embedding of $t$}.
  \begin{itemize}
    \item How to embed $(h,r)$? How to define score $f_r(h,t)$?
    \item Score $f_r$ is high if $(h,r,t)$ exists, else $f_r$ is low.
  \end{itemize}
\end{keyidea}

\begin{figure}[htbp]
  \centering
  \begin{tikzpicture}[scale=0.85]
    \node[lec10gen] (sf)  at (3.0,5.5) {};
    \node[lec10ent] (sk)  at (4.2,2.9) {};
    \node[lec10ent] (ap)  at (2.7,0.5) {};
    \node[lec10gen] (fan) at (7.3,3.9) {};
    \node[lec10ent] (jkr) at (7.5,2.0) {};
    \node[lec10gen] (tra) at (6.4,0.5) {};
    \node[lec10ent] (la)  at (9.1,0.2) {};
    \node[lec10ent] (csl) at (10.8,2.0) {};
    \node[lec10ent] (tol) at (8.4,6.2) {};
    \node[lec10name, ForestGreen, left=1pt of sf] {Science Fiction};
    \node[lec10name, RoyalBlue, left=1pt of sk] {Stephen King};
    \node[lec10name, RoyalBlue, below=1pt of ap] {Alan Poe};
    \node[lec10name, ForestGreen, below=1pt of fan] {Fantasy};
    \node[lec10name, RoyalBlue, above right=0pt and 1pt of jkr] {J.K. Rowling};
    \node[lec10name, ForestGreen, below=1pt of tra] {Tragicomedy};
    \node[lec10name, RoyalBlue, right=1pt of la] {Lloyd Alexander};
    \node[lec10name, RoyalBlue, below=1pt of csl] {C.S. Lewis};
    \node[lec10name, RoyalBlue, right=1pt of tol] {J.R.R Tolkien};
    \draw[lec10gr]  (sk)  -- node[lec10lbl] {genre} (sf);
    \draw[lec10gr]  (sk)  -- node[lec10lbl, pos=0.22] {genre} (fan);
    \draw[lec10inf] (sk)  -- node[lec10lbl] {Influence} (ap);
    \draw[lec10inf] (sk)  -- node[lec10lbl, pos=0.45] {Influence} (jkr);
    \draw[lec10inf] (jkr) -- node[lec10lbl] {Influence} (csl);
    \draw[lec10inf] (jkr) -- node[lec10lbl, pos=0.55] {Influence} (la);
    \draw[lec10gr]  (jkr) -- node[lec10lbl] {genre} (tra);
    \draw[lec10gr]  (csl) -- node[lec10lbl, pos=0.62] {genre} (fan);
    \draw[lec10inf] (csl) -- node[lec10lbl] {Influence} (tol);
    \draw[lec10gr]  (tol) -- node[lec10lbl] {genre} (fan);
    \draw[thick, densely dashed, magenta, -{Stealth[length=1.8mm]}] (jkr) -- (sf);
    \node[font=\scriptsize\sffamily, magenta, align=center] at (6.0,5.4)
      {missing relation:\\genre};
  \end{tikzpicture}
  \caption{KG completion: predict the missing tail ``Science Fiction'' for
  (``J.K. Rowling'', ``genre'') (slide~14).}
  \label{fig:lec10-completion}
\end{figure}

\subsection{KG embedding models}\label{sec:lec10-models}
\slides{16--17}

\begin{itemize}
  \item \textbf{Many KG embedding models} (timeline 2011--2020; three colour
    groups on the slide):
    \begin{itemize}
      \item \textcolor{RoyalBlue}{blue}: SE (Bordes et al., 2011), UM (Bordes et
        al., 2012), TransE (Bordes et al., 2013), TransH (Wang et al., 2014),
        KG2E (He et al., 2015), TransR (Lin et al., 2015), Poincar\'e (Nickel et
        al., 2017), TorusE (He et al., 2018), RotatE (Sun et al., 2019), BoxE
        (Abboud et al., 2020);
      \item \textcolor{ForestGreen}{green}: RESCAL (Nickel et al., 2011), DistMult
        (Yang et al., 2015), ComplEx (Trouillon et al., labelled 2014, placed at
        2016), TuckER (Balazevic et al., 2019);
      \item \textcolor{orange}{orange}: NTN (Socher et al., 2013), SME (Bordes et
        al., 2014), ConvE (Dettmers et al., 2018), r-GCN (Schlichtkrull et al.,
        2018).
    \end{itemize}
  \item \textbf{Today:} different KG embedding models
    (\textbf{shallow/transductive embs}). Different models are \dots
    \begin{itemize}
      \item \dots based on different \textbf{geometric intuitions};
      \item \dots capture different types of relations (have different
        \textbf{expressivity}) -- the overview table of slide~17 is
        \cref{tab:lec10-patterns} (without the RotatE row).
    \end{itemize}
\end{itemize}

% ---------------------------------------------------------------------
\section{TransE}\label{sec:lec10-transe}

\subsection{Translation intuition}\label{sec:lec10-transe-def}
\slides{19--20}

\textbf{TransE} (Bordes et al., \emph{Translating embeddings for modeling
multi-relational data}, NeurIPS 2013).
\begin{itemize}
  \item \textbf{Intuition: translation.} For a triple $(h,r,t)$, let
    $\vh,\vr,\vt\in\R^k$ be embedding vectors (embedding vectors appear in
    boldface).
  \item \textbf{TransE:} $\vh+\vr\approx\vt$ if the given link exists, else
    $\vh+\vr\neq\vt$. \textbf{Entity scoring function:}
\end{itemize}
\begin{equation}\label{eq:lec10-transe}
  f_r(h,t) = -\norm{\vh+\vr-\vt}
\end{equation}
\begin{itemize}
  \item Example: Obama $+$ Nationality $\approx$ U.S.A
    (\cref{fig:lec10-transe}a).
  \item \textbf{TransE idea -- entity embeddings:} king $\xrightarrow{female}$
    queen and man $\xrightarrow{female}$ woman; man $\xrightarrow{royal}$ king and
    woman $\xrightarrow{royal}$ queen -- a parallelogram
    (\cref{fig:lec10-transe}b).
\end{itemize}

\subsection{How to learn TransE}\label{sec:lec10-transe-learn}
\slides{21}

\begin{itemize}
  \item \textbf{Input:} training set $S=\{(h,r,t)\}$, entities and rel.\ sets $E$
    and $R$, margin $\gamma$, embeddings dim.\ $k$ (\cref{alg:lec10-transe}).
  \item Initialize relations $r$ and entities $e$ \textbf{uniformly}, then
    \textbf{normalize}; $\gamma$ is the \textbf{margin}.
  \item For each training triple, sample a \textbf{corrupted triplet} $(h',r,t)$
    that \textbf{does not appear in the KG}.
  \item Update embeddings w.r.t.\ the \term{contrastive loss} ($d$ represents
    distance = negative of score):
\end{itemize}
\begin{equation}\label{eq:lec10-loss}
  \sum_{((h,r,t),(h',r,t'))\in T_{batch}}
  \big[\gamma + \underbrace{d(\vh+\vr,\vt)}_{\text{positive sample}}
             - \underbrace{d(\vh'+\vr,\vt')}_{\text{negative sample}}\big]_+
\end{equation}
\begin{itemize}
  \item[] Contrastive loss: favors \textbf{lower distance (or higher score) for
    valid triplets}, \textbf{high distance (or lower score) for corrupted ones}.
\end{itemize}

\begin{algorithm}[htbp]
  \caption{Learning TransE (slide~21)}\label{alg:lec10-transe}
  \begin{algorithmic}[1]
    \Statex \textbf{input} Training set $S=\{(h,r,t)\}$, entities and rel.\ sets
      $E$ and $R$, margin $\gamma$, embeddings dim.\ $k$.
    \State \textbf{initialize} $\vr\gets\mathrm{uniform}(-\frac{6}{\sqrt{k}},
      \frac{6}{\sqrt{k}})$ for each $r\in R$
    \State \phantom{\textbf{initialize}} $\vr\gets\vr/\norm{\vr}$ for each $r\in R$
    \State \phantom{\textbf{initialize}} $\ve\gets\mathrm{uniform}(-\frac{6}{\sqrt{k}},
      \frac{6}{\sqrt{k}})$ for each entity $e\in E$
    \Loop
      \State $\ve\gets\ve/\norm{\ve}$ for each entity $e\in E$
      \State $S_{batch}\gets\mathrm{sample}(S,b)$ \Comment{sample a minibatch of size $b$}
      \State $T_{batch}\gets\emptyset$ \Comment{initialize the set of pairs of triplets}
      \For{$(h,r,t)\in S_{batch}$}
        \State $(h',r,t')\gets\mathrm{sample}(S'_{(h,r,t)})$ \Comment{sample a corrupted triplet}
        \State $T_{batch}\gets T_{batch}\cup\{((h,r,t),(h',r,t'))\}$
      \EndFor
      \State Update embeddings w.r.t.\ the loss \eqref{eq:lec10-loss}
    \EndLoop
  \end{algorithmic}
\end{algorithm}

\subsection{Relation patterns}\label{sec:lec10-patterns}
\slides{22--23}

\begin{itemize}
  \item \textbf{Relations in a heterogeneous KG have different properties.}
    Example:
    \begin{itemize}
      \item \textbf{Symmetry:} if the edge $(h,\text{``Roommate''},t)$ exists
        in KG, then the edge\\ $(t,\text{``Roommate''},h)$ should also exist.
      \item \textbf{Inverse relation:} if the edge $(h,\text{``Advisor''},t)$
        exists in KG, then the edge\\ $(t,\text{``Advisee''},h)$ should also
        exist.
    \end{itemize}
  \item Can we \textbf{categorize} these relation patterns? \textbf{Are KG
    embedding methods (e.g., TransE) expressive enough to model these
    patterns?}
\end{itemize}

\begin{definition}[Four relation patterns]\label{def:lec10-patterns}
  \begin{align}
    \text{\textbf{Symmetric (Antisymmetric):}}\quad
      & r(h,t)\Rightarrow r(t,h)\ \ \big(r(h,t)\Rightarrow\neg r(t,h)\big)
        \quad\forall h,t \label{eq:lec10-sym}\\
    \text{\textbf{Inverse:}}\quad
      & r_2(h,t)\Rightarrow r_1(t,h) \label{eq:lec10-inv}\\
    \text{\textbf{Composition (Transitive):}}\quad
      & r_1(x,y)\wedge r_2(y,z)\Rightarrow r_3(x,z)\quad\forall x,y,z
        \label{eq:lec10-comp}\\
    \text{\textbf{1-to-N:}}\quad
      & r(h,t_1), r(h,t_2),\dots,r(h,t_n)\ \text{are all True.}
        \label{eq:lec10-1toN}
  \end{align}
  \textbf{Examples:} symmetric -- Family, Roommate; antisymmetric -- Hypernym (a
  word with a broader meaning: poodle vs.\ dog); inverse -- (Advisor, Advisee);
  composition -- my mother's husband is my father; 1-to-N -- $r$ is
  ``StudentsOf''.
\end{definition}

\subsection{Relation patterns in TransE}\label{sec:lec10-transe-patterns}
\slides{24--29}

\begin{itemize}
  \item \textbf{Antisymmetric \cmark:} $\vh+\vr=\vt$, but $\vt+\vr\neq\vh$
    (\cref{fig:lec10-transe}c).
  \item \textbf{Inverse \cmark:} $\vh+\vr_2=\vt$, we can set $\vr_1=-\vr_2$
    (\cref{fig:lec10-transe}d).
  \item \textbf{Composition \cmark:} $\vr_3=\vr_1+\vr_2$
    (\cref{fig:lec10-transe}e).
  \item \textbf{Symmetric \xmark} -- \textbf{only if $\vr=0$, $\vh=\vt$}: for all
    $h,t$ that satisfy $r(h,t)$, $r(t,h)$ is also True, which means
    $\norm{\vh+\vr-\vt}=0$ and $\norm{\vt+\vr-\vh}=0$. Then $\vr=0$ and
    $\vh=\vt$; however $h$ and $t$ are two different entities and should be
    mapped to different locations.
  \item \textbf{1-to-N \xmark:} $(h,r,t_1)$ and $(h,r,t_2)$ both exist in the
    knowledge graph (e.g.\ $r$ is ``StudentsOf''); $\vt_1$ and $\vt_2$ will map to
    the same vector, although they are different entities
    (\cref{fig:lec10-transe}f):
    \[
      \vt_1=\vh+\vr=\vt_2,\qquad \vt_1\neq\vt_2\qquad\text{contradictory!}
    \]
  \item \textbf{What we learned so far} (slide~29): the TransE row of
    \cref{tab:lec10-patterns}.
\end{itemize}

\begin{figure}[htbp]
  \centering
  \begin{tikzpicture}[scale=0.85]
    % (a) translation
    \begin{scope}
      \draw[lec10ax] (0,2.2) -- (0,0) -- (3.2,0);
      \node[lec10pt] (h) at (0.6,0.7) {}; \node[lec10pt] (t) at (2.5,1.2) {};
      \draw[lec10vec, red!75!black] (h) -- node[above, font=\small] {$\vr$} (t);
      \node[left, blue, font=\small] at (h) {$\vh$};
      \node[above, ForestGreen, font=\small] at (t) {$\vt$};
      \node[lec10cap] at (2.1,1.95) {Obama $\xrightarrow{\text{Nationality}}$ U.S.A};
      \node[lec10cap] at (1.6,-0.4) {(a) $\vh+\vr\approx\vt$ (sl.~19)};
    \end{scope}
    % (b) king / queen
    \begin{scope}[xshift=4.3cm]
      \node[font=\scriptsize] (k) at (0.5,1.8) {king};
      \node[font=\scriptsize] (q) at (2.7,1.8) {queen};
      \node[font=\scriptsize] (m) at (0.1,0.3) {man};
      \node[font=\scriptsize] (w) at (2.3,0.3) {woman};
      \draw[lec10vec, red!75!black] (k) -- node[above, font=\tiny\itshape] {female} (q);
      \draw[lec10vec, red!75!black] (m) -- node[below, font=\tiny\itshape] {female} (w);
      \draw[lec10vec, blue!70] (m) -- node[left, font=\tiny\itshape] {royal} (k);
      \draw[lec10vec, blue!70] (w) -- node[right, font=\tiny\itshape] {royal} (q);
      \node[lec10cap] at (1.4,-0.4) {(b) entity embeddings (sl.~20)};
    \end{scope}
    % (c) antisymmetric
    \begin{scope}[xshift=8.6cm]
      \draw[lec10ax] (0,2.2) -- (0,0) -- (3.2,0);
      \node[lec10pt] (h) at (0.5,0.6) {}; \node[lec10pt] (t) at (1.5,0.9) {};
      \draw[lec10vec, red!75!black] (h) -- node[above, font=\small] {$\vr$} (t);
      \draw[lec10vec, red!75!black] (t) -- node[above, font=\small] {$\vr$} (2.6,1.25);
      \node[below, font=\small] at (h) {$\vh$};
      \node[below, font=\small] at (t) {$\vt$};
      \node[lec10cap] at (1.6,-0.4) {(c) antisymmetric \cmark\ (sl.~24)};
    \end{scope}
    % (d) inverse
    \begin{scope}[yshift=-3.3cm]
      \draw[lec10ax] (0,2.2) -- (0,0) -- (3.2,0);
      \node[lec10pt] (h) at (0.6,0.8) {}; \node[lec10pt] (t) at (2.4,1.2) {};
      \draw[lec10vec, red!75!black] ([yshift=1.5pt]h.north east) --
        node[above, font=\small] {$\vr_1$} ([yshift=1.5pt]t.north west);
      \draw[lec10vec, ForestGreen] ([yshift=-1.5pt]t.south west) --
        node[below, font=\small] {$\vr_2$} ([yshift=-1.5pt]h.south east);
      \node[left, font=\small] at (h) {$\vh$};
      \node[right, font=\small] at (t) {$\vt$};
      \node[lec10cap] at (1.6,-0.4) {(d) inverse \cmark\ (sl.~25)};
    \end{scope}
    % (e) composition
    \begin{scope}[xshift=4.3cm, yshift=-3.3cm]
      \draw[lec10ax] (0,2.2) -- (0,0) -- (3.2,0);
      \node[lec10pt] (x) at (0.5,0.5) {}; \node[lec10pt] (y) at (1.4,1.5) {};
      \node[lec10pt] (z) at (2.7,0.7) {};
      \draw[lec10vec, red!75!black] (x) -- node[above left=-2pt, font=\small] {$\vr_1$} (y);
      \draw[lec10vec, orange] (y) -- node[above right=-2pt, font=\small] {$\vr_2$} (z);
      \draw[lec10vec, blue] (x) -- node[below, font=\small] {$\vr_3$} (z);
      \node[left, font=\small] at (x) {$\vx$};
      \node[above, font=\small] at (y) {$\vy$};
      \node[right, font=\small] at (z) {$\vz$};
      \node[lec10cap] at (1.6,-0.4) {(e) composition \cmark\ (sl.~26)};
    \end{scope}
    % (f) 1-to-N
    \begin{scope}[xshift=8.6cm, yshift=-3.3cm]
      \draw[lec10ax] (0,2.2) -- (0,0) -- (3.2,0);
      \node[lec10pt] (h) at (2.3,0.4) {}; \node[lec10pt] (t1) at (1.0,1.8) {};
      \node[lec10pt] (t2) at (1.1,1.1) {};
      \draw[lec10vec] (h) -- node[right, font=\small] {$\vr$} (t1);
      \draw[lec10vec] (h) -- node[below, font=\small] {$\vr$} (t2);
      \node[below right=-1pt, font=\small] at (h) {$\vh$};
      \node[left, red!75!black, font=\small] at (t1) {$\vt_1$};
      \node[left, ForestGreen, font=\small] at (t2) {$\vt_2$};
      \node[lec10cap] at (1.6,-0.4) {(f) 1-to-N \xmark\ (sl.~28)};
    \end{scope}
  \end{tikzpicture}
  \caption{TransE geometry (slides~19--28).}
  \label{fig:lec10-transe}
\end{figure}

% ---------------------------------------------------------------------
\section{TransR}\label{sec:lec10-transr}

\subsection{Translation in relation-specific space}\label{sec:lec10-transr-def}
\slides{31--32}

\textbf{TransR} (Lin et al., \emph{Learning entity and relation embeddings for
knowledge graph completion}, AAAI 2015).
\begin{itemize}
  \item TransE models translation of any relation in the \textbf{same} embedding
    space. \textbf{Can we design a new space for each relation and do
    translation in relation-specific space?}
  \item \textbf{TransR:} model \textbf{entities} as vectors in the entity space
    $\R^d$ and model each \textbf{relation} as vector in relation space
    $\vr\in\R^k$ with $\mM_r\in\R^{k\times d}$ as the \term{projection matrix}.
    Use $\mM_r$ to \textbf{project} from entity space $\R^d$ to relation space
    $\R^k$! (\cref{fig:lec10-transr}a)
\end{itemize}
\begin{equation}\label{eq:lec10-transr}
  \vh_\perp=\mM_r\vh,\quad \vt_\perp=\mM_r\vt,\qquad
  f_r(h,t) = -\norm{\vh_\perp+\vr-\vt_\perp}
\end{equation}

\subsection{Relation patterns in TransR}\label{sec:lec10-transr-patterns}
\slides{33--36}

\begin{itemize}
  \item \textbf{Symmetric \cmark:}
    $\vr=0,\ \vh_\perp=\mM_r\vh=\mM_r\vt=\vt_\perp$. We can map $h$ and $t$ to the
    same location on the space of relation $r$; then $r=0$; $h$ and $t$ are still
    different in the entity space. Note different symmetric relations may have
    different $\mM_r$.
  \item \textbf{Antisymmetric \cmark:} $\vr\neq0$, $\mM_r\vh+\vr=\mM_r\vt$; then
    $\mM_r\vt+\vr\neq\mM_r\vh$.
  \item \textbf{1-to-N \cmark:} we can learn $\mM_r$ so that
    $\vt_\perp=\mM_r\vt_1=\mM_r\vt_2$. Note that $\vt_1$ does not need to be equal
    to $\vt_2$! (\cref{fig:lec10-transr}b)
  \item \textbf{Inverse \cmark:} $\vr_2=-\vr_1$, $\mM_{r_1}=\mM_{r_2}$. Then
    $\mM_{r_1}\vt+\vr_1=\mM_{r_1}\vh$ and $\mM_{r_2}\vh+\vr_2=\mM_{r_2}\vt$.
\end{itemize}

\begin{figure}[htbp]
  \centering
  \begin{tikzpicture}[scale=0.85]
    % (a) projection
    \begin{scope}
      \draw[lec10ax] (0,2.6) -- (0,0) -- (3.0,0);
      \draw[lec10ax] (3.4,2.6) -- (3.4,0) -- (6.6,0);
      \node[font=\tiny\sffamily\bfseries, ForestGreen] at (1.5,2.8) {Space of entities: $\R^d$};
      \node[font=\tiny\sffamily\bfseries, blue] at (5.0,2.8) {Space of relation $r$: $\R^k$};
      \node[lec10pt] (t) at (0.5,2.0) {}; \node[lec10pt] (h) at (1.9,1.1) {};
      \node[lec10pt] (hp) at (4.4,0.8) {}; \node[lec10pt] (tp) at (5.9,1.15) {};
      \draw[lec10proj] (t) -- (tp); \draw[lec10proj] (h) -- (hp);
      \draw[lec10vec] (hp) -- node[below, font=\small] {$\vr$} (tp);
      \node[below, font=\small] at (t) {$\vt$};
      \node[below, font=\small] at (h) {$\vh$};
      \node[below, font=\small] at (hp) {$\vh_\perp$};
      \node[above, font=\small] at (tp) {$\vt_\perp$};
      \node[blue, font=\small] at (3.0,1.85) {$\mM_r$};
      \node[lec10cap] at (3.3,-0.4) {(a) TransR projection (sl.~32)};
    \end{scope}
    % (b) 1-to-N
    \begin{scope}[xshift=7.4cm]
      \draw[lec10ax] (0,2.6) -- (0,0) -- (3.0,0);
      \draw[lec10ax] (3.4,2.6) -- (3.4,0) -- (6.6,0);
      \node[lec10pt] (t1) at (0.8,2.1) {}; \node[lec10pt] (t2) at (0.9,1.4) {};
      \node[lec10pt] (h) at (1.4,0.6) {};
      \node[lec10pt] (hp) at (4.4,1.0) {}; \node[lec10pt] (tp) at (5.9,1.35) {};
      \draw[lec10proj] (t1) -- (tp); \draw[lec10proj] (t2) -- (tp);
      \draw[lec10proj] (h) -- (hp);
      \draw[lec10vec] (hp) -- node[below, font=\small] {$\vr$} (tp);
      \node[left, red!75!black, font=\small] at (t1) {$\vt_1$};
      \node[left, ForestGreen, font=\small] at (t2) {$\vt_2$};
      \node[below, font=\small] at (h) {$\vh$};
      \node[below, font=\small] at (hp) {$\vh_\perp$};
      \node[above right=-2pt, font=\small] at (tp) {$\vt_\perp$};
      \node[blue, font=\small] at (2.6,2.05) {$\mM_r$};
      \node[lec10cap] at (3.3,-0.4) {(b) 1-to-N: $\mM_r\vt_1=\mM_r\vt_2$ (sl.~35)};
    \end{scope}
  \end{tikzpicture}
  \caption{TransR: entities in $\R^d$ are projected by $\mM_r$ into the space of
  relation $r$, $\R^k$, where $\vh_\perp+\vr\approx\vt_\perp$ (slides~32, 35).}
  \label{fig:lec10-transr}
\end{figure}

\subsection{Composition relations in TransR}\label{sec:lec10-transr-comp}
\slides{37--41}

\begin{itemize}
  \item \textbf{Composition \cmark.} \textbf{High-level intuition:} TransR models
    a triple with linear functions. \textbf{Linear functions are chainable!} If
    $f(x)$ and $g(x)$ are linear, then $f(g(x))$ is also linear: let
    $f(x)=a\cdot x+b$, $g(x)=c\cdot x+d$; then $f(g(x))=a(c\cdot x+d)+b$.
  \item \textbf{Background -- Def: kernel space of a matrix $\mM$:}
    $\vh\in\Ker(\mM)$, then $\mM\cdot\vh=\mathbf{0}$ (figure: the region
    $\Ker(\mM)$ is mapped by $\mM$ to $\mathbf{0}$).
  \item \textbf{Argument.} Assume $\mM_{r_1}\vg_1=\vr_1$ and
    $\mM_{r_2}\vg_2=\vr_2$.
\end{itemize}
\begin{align}
  \text{For } r_1(x,y):\ \ & r_1(x,y)\text{ exists}\Rightarrow
    \mM_{r_1}\vx+\vr_1=\mM_{r_1}\vy\Rightarrow\mM_{r_1}(\vy-\vx)=\vr_1 \notag\\
  & \vy-\vx\in\vg_1+\Ker(\mM_{r_1})\Rightarrow\vy\in\vx+\vg_1+\Ker(\mM_{r_1})\notag\\
  \text{Same for } r_2(y,z):\ \ & r_2(y,z)\text{ exists}\Rightarrow
    \mM_{r_2}\vy+\vr_2=\mM_{r_2}\vz\Rightarrow\notag\\
  & \vz-\vy\in\vg_2+\Ker(\mM_{r_2})\Rightarrow\vz\in\vy+\vg_2+\Ker(\mM_{r_2})\notag\\
  \text{Then, we have}\ \ & \vz\in\vx+\vg_1+\vg_2+\Ker(\mM_{r_1})+\Ker(\mM_{r_2})
    \label{eq:lec10-kernel}
\end{align}
\begin{itemize}
  \item \textbf{Construct $\mM_{r_3}$} s.t.\
    $\Ker(\mM_{r_3})=\Ker(\mM_{r_1})+\Ker(\mM_{r_2})$. Since
    $\dim\big(\Ker(\mM_{r_3})\big)\ge\dim\big(\Ker(\mM_{r_1})\big)$ and
    $\mM_{r_3}$ has the same shape as $\mM_{r_1}$, we know $\mM_{r_3}$ exists!
  \item \textbf{Set} $\vr_3=\mM_{r_3}(\vg_1+\vg_2)$. \textbf{We have}
    $\mM_{r_3}\vx+\vr_3=\mM_{r_3}\vz$.
  \item \textbf{What we learned so far} (slide~41): TransR captures all five
    patterns (\cref{tab:lec10-patterns}).
\end{itemize}

% ---------------------------------------------------------------------
\section{DistMult}\label{sec:lec10-distmult}

\subsection{Bilinear modeling}\label{sec:lec10-bilinear}
\slides{43--45}

\textbf{DistMult} (Yang et al., \emph{Embedding Entities and Relations for
Learning and Inference in Knowledge Bases}, ICLR 2015).
\begin{itemize}
  \item \textbf{So far:} the scoring function $f_r(h,t)$ is the \textbf{negative
    of L1 / L2 distance} in TransE and TransR.
  \item \textbf{Idea: use bilinear modeling.} Score function:
\end{itemize}
\begin{equation}\label{eq:lec10-bilinear}
  f_r(h,t)=h\cdot A\cdot t,\qquad \vh,\vt\in\R^k,\ \mA\in\R^{k\times k}
\end{equation}
\begin{itemize}
  \item \textbf{Problem:} too general and prone to overfitting -- matrix $A$ is
    too expressive. \textbf{Fix: limit $A$ to be diagonal} $\Rightarrow$ this is
    called \term{DistMult}.
  \item \textbf{DistMult:} entities and relations are vectors in $\R^k$. Score
    function (elementwise product, then sum):
\end{itemize}
\begin{equation}\label{eq:lec10-distmult}
  f_r(h,t)=\langle\vh,\vr,\vt\rangle=\sum_i\vh_i\cdot\vr_i\cdot\vt_i,
  \qquad \vh,\vr,\vt\in\R^k
\end{equation}
\begin{itemize}
  \item \textbf{Intuition of the score function:} can be viewed as a
    \textbf{cosine similarity} between $\vh\cdot\vr$ and $\vt$, where
    $\vh\cdot\vr$ is defined as $[\vh\cdot\vr]_i=\vh_i\cdot\vr_i$
    (\textbf{Hadamard product}):
\end{itemize}
\begin{equation}\label{eq:lec10-cos}
  \cos(\vh\cdot\vr,\vt)=\frac{\langle\vh,\vr,\vt\rangle}{\norm{\vh\cdot\vr}\norm{\vt}},
  \qquad
  \langle\vh,\vr,\vt\rangle=\cos(\vh\cdot\vr,\vt)\norm{\vh\cdot\vr}\norm{\vt}
\end{equation}
\begin{itemize}
  \item \textbf{Example:} $f_r(h,t_1)<0$, $f_r(h,t_2)>0$ ($\vt_2$ on the
    $\vh\cdot\vr$ side of the dashed line orthogonal to $\vh\cdot\vr$, $\vt_1$ on
    the other side; \cref{fig:lec10-distmult}a).
\end{itemize}

\subsection{Relation patterns in DistMult}\label{sec:lec10-distmult-patterns}
\slides{46--51}

\begin{itemize}
  \item \textbf{1-to-N \cmark:}
    $\langle\vh,\vr,\vt_1\rangle=\langle\vh,\vr,\vt_2\rangle$ ($\vt_1,\vt_2$ on
    the same dashed line orthogonal to $\vh\cdot\vr$,
    \cref{fig:lec10-distmult}b).
  \item \textbf{Symmetric \cmark} (DistMult can \textbf{naturally} model them),
    due to the commutative property of multiplication:
    \[
      f_r(h,t)=\langle\vh,\vr,\vt\rangle=\sum_i\vh_i\cdot\vr_i\cdot\vt_i
      =\langle\vt,\vr,\vh\rangle=f_r(t,h)
    \]
  \item \textbf{Antisymmetric \xmark:}
    $f_r(h,t)=\langle\vh,\vr,\vt\rangle=\langle\vt,\vr,\vh\rangle=f_r(t,h)$ --
    $r(h,t)$ and $r(t,h)$ \textbf{always have same score!} DistMult cannot
    differentiate between head entity and tail entity: all relations are
    modelled as symmetric regardless, i.e.\ even antisymmetric relations will be
    represented as symmetric.
  \item \textbf{Inverse \xmark:} assume DistMult does model inverse relations:
    $f_{r_2}(h,t)=\langle\vh,\vr_2,\vt\rangle=\langle\vt,\vr_1,\vh\rangle=f_{r_1}(t,h)$.
    For example, $\vr_2=\vr_1$ solves this (there also exist solutions
    $\vr_2\neq\vr_1$). \textbf{But semantically this does not make sense: the
    embedding of ``Advisor'' relation should not be the same as ``Advisee''
    relation.}
  \item \textbf{Composition \xmark.} Intuition: because dot product is
    commutative ($a\cdot b=b\cdot a$), DistMult does not distinguish between head
    and tail entities, so it cannot model composition.
  \item \textbf{What we learned so far} (slide~51): the DistMult row of
    \cref{tab:lec10-patterns}.
\end{itemize}

\begin{figure}[htbp]
  \centering
  \begin{tikzpicture}[scale=0.85]
    % (a) sign of the score
    \begin{scope}
      \coordinate (o) at (0,0);
      \draw[lec10proj, thick] (0.33,-1.43) -- (-0.4,1.73);
      \draw[lec10vec] (o) -- (2.8,0.65) node[right, font=\small] {$\vh\cdot\vr$};
      \draw[lec10vec, ForestGreen] (o) -- (2.0,1.4) node[above, font=\small] {$\vt_2$};
      \draw[lec10vec, red!75!black] (o) -- (-0.7,-0.85) node[below left=-2pt, font=\small] {$\vt_1$};
      \node[font=\scriptsize, align=left] at (2.3,-0.9) {$f_r(h,t_1)<0$\\ $f_r(h,t_2)>0$};
      \node[lec10cap] at (1.0,-1.9) {(a) sign of $\langle\vh,\vr,\vt\rangle$ (sl.~45)};
    \end{scope}
    % (b) 1-to-N
    \begin{scope}[xshift=6.0cm, yshift=-1.4cm]
      \draw[lec10ax] (0,3.4) -- (0,0) -- (3.4,0);
      \draw[lec10proj] (2.24,0.18) -- (1.75,2.9);
      \draw[lec10vec] (0,0) -- (2.2,0.4) node[right, font=\small] {$\vh\cdot\vr$};
      \draw[lec10vec, ForestGreen] (0,0) -- (2.02,1.39) node[right, font=\small] {$\vt_2$};
      \draw[lec10vec, red!75!black] (0,0) -- (1.9,2.05) node[right, font=\small] {$\vt_1$};
      \node[lec10cap] at (1.7,-0.5) {(b) 1-to-N:
        $\langle\vh,\vr,\vt_1\rangle=\langle\vh,\vr,\vt_2\rangle$ (sl.~46)};
    \end{scope}
  \end{tikzpicture}
  \caption{DistMult geometry (slides~45--46).}
  \label{fig:lec10-distmult}
\end{figure}

% ---------------------------------------------------------------------
\section{ComplEx}\label{sec:lec10-complex}

\subsection{Complex embeddings and the score function}\label{sec:lec10-complex-def}
\slides{53--55}

\textbf{ComplEx} (Trouillon et al., \emph{Complex Embeddings for Simple Link
Prediction}, ICML 2016).
\begin{itemize}
  \item Based on DistMult, \textbf{ComplEx} embeds entities and relations in
    \textbf{complex vector space}: vectors in $\C^k$.
  \item $\vu\in\C^k$: $\vu=\va+\vb i$ with $\va\in\R^k$, $\vb\in\R^k$;
    $\bar\vu=\va-\vb i$ is called a \term{conjugate} (figure: $\bar\vu$ is the
    mirror image of $\vu$ in the plane with axes $\cRe(\vu)$, $\cIm(\vu)$).
  \item \textbf{Complex multiplication:} $(a+ib)(c+id)=(ac-bd)+i(ad+bc)$.
    Example: $(3+4i)\cdot(2+3i)=(3+4i)\cdot 2+(3+4i)\cdot 3i=-6+17i$ (scale
    ($\ast2$) plus scale/rotate ($\ast3i$)).
  \item \textbf{Score function} (inner product of $\vh$, $\vr$, $\bar\vt$, then
    sum):
\end{itemize}
\begin{equation}\label{eq:lec10-complex}
  f_r(h,t)=\cRe\Big(\sum_i\vh_i\cdot\vr_i\cdot\bar\vt_i\Big)
\end{equation}
Expanded (slides~54--55):
\begin{align}
  f_r(h,t)&=\cRe\Big(\sum_i\vh_i\cdot\vr_i\cdot\bar\vt_i\Big)
    =\sum_i\cRe(\vh_i\cdot\vr_i\cdot\bar\vt_i)\notag\\
  &=\sum_i\cRe\Big(\big(\cRe(\vh_i)+i\cIm(\vh_i)\big)\cdot
    \big(\cRe(\vr_i)+i\cIm(\vr_i)\big)\cdot\big(\cRe(\vt_i)-i\cIm(\vt_i)\big)\Big)\notag\\
  &=\sum_i\Big[\cRe(\vh_i)\cRe(\vr_i)\cRe(\vt_i)+\cRe(\vh_i)\cIm(\vr_i)\cIm(\vt_i)\notag\\
  &\qquad\quad +\cIm(\vh_i)\cRe(\vr_i)\cIm(\vt_i)-\cIm(\vh_i)\cIm(\vr_i)\cRe(\vt_i)\Big]\notag\\
  &=\langle\cRe(\vh_i),\cRe(\vr_i),\cRe(\vt_i)\rangle
    +\langle\cRe(\vh_i),\cIm(\vr_i),\cIm(\vt_i)\rangle\notag\\
  &\quad+\langle\cIm(\vh_i),\cRe(\vr_i),\cIm(\vt_i)\rangle
    -\langle\cIm(\vh_i),\cIm(\vr_i),\cRe(\vt_i)\rangle
  \label{eq:lec10-complex-expand}
\end{align}

\subsection{Relation patterns in ComplEx}\label{sec:lec10-complex-patterns}
\slides{56--59}

\begin{itemize}
  \item \textbf{Antisymmetric \cmark:} the model is expressive enough to learn
    \begin{itemize}
      \item \textbf{high} $f_r(h,t)=\cRe\big(\sum_i\vh_i\cdot\vr_i\cdot\bar\vt_i\big)$,
      \item \textbf{low} $f_r(t,h)=\cRe\big(\sum_i\vt_i\cdot\vr_i\cdot\bar\vh_i\big)$,
    \end{itemize}
    due to the \textbf{asymmetric modeling using complex conjugate}.
  \item \textbf{Symmetric \cmark:} when $\cIm(\vr)=0$, we have
\end{itemize}
\begin{align}
  f_r(h,t)&=\cRe\Big(\sum_i\vh_i\cdot\vr_i\cdot\bar\vt_i\Big)
    =\sum_i\cRe(\vr_i\cdot\vh_i\cdot\bar\vt_i)
    =\sum_i\vr_i\cdot\cRe(\vh_i\cdot\bar\vt_i)\notag\\
  &=\sum_i\vr_i\cdot\cRe(\bar\vh_i\cdot\vt_i)
    =\sum_i\cRe(\vr_i\cdot\bar\vh_i\cdot\vt_i)=f_r(t,h)
  \label{eq:lec10-complex-sym}
\end{align}
\begin{itemize}
  \item \textbf{Inverse \cmark:} $\vr_1=\bar\vr_2$. The complex conjugate of
    $\vr_2=\argmax_{\vr}\cRe(\langle\vh,\vr,\bar\vt\rangle)$ is exactly
    $\vr_1=\argmax_{\vr}\cRe(\langle\vt,\vr,\bar\vh\rangle)$.
  \item \textbf{Composition and 1-to-N:} ComplEx shares the same property with
    DistMult -- \textbf{cannot model composition relations} \xmark; \textbf{can
    model 1-to-N relations} \cmark.
\end{itemize}

% ---------------------------------------------------------------------
\section{Comparing KG embedding models}\label{sec:lec10-compare}

\subsection{Expressiveness of the models}\label{sec:lec10-table}
\slides{60}

\begin{table}[htbp]
  \centering\small
  \caption{KG completion models: score, embedding space, and the relation
  patterns each can model (slides~17, 29, 41, 51, 60).}
  \label{tab:lec10-patterns}
  \renewcommand{\arraystretch}{1.3}
  \begin{tabular}{@{}llccccccc@{}}
    \toprule
    Model & Score & Embedding & Sym. & Antisym. & Inv. & Compos. & 1-to-N\\
    \midrule
    TransE & $-\norm{\vh+\vr-\vt}$ & $\vh,\vt,\vr\in\R^k$
      & \xmark & \cmark & \cmark & \cmark & \xmark\\
    TransR & $-\norm{\mM_r\vh+\vr-\mM_r\vt}$
      & $\begin{array}{@{}c@{}}\vh,\vt\in\R^d,\ \vr\in\R^k,\\ \mM_r\in\R^{k\times d}\end{array}$
      & \cmark & \cmark & \cmark & \cmark & \cmark\\
    DistMult & $\langle\vh,\vr,\vt\rangle$ & $\vh,\vt,\vr\in\R^k$
      & \cmark & \xmark & \xmark & \xmark & \cmark\\
    ComplEx & $\cRe(\langle\vh,\vr,\bar\vt\rangle)$ & $\vh,\vt,\vr\in\C^k$
      & \cmark & \cmark & \cmark & \xmark & \cmark\\
    RotatE & $-\norm{\vh\circ\vr-\vt}$ & $\vh,\vt,\vr\in\C^k$
      & \cmark & \cmark & \cmark & \cmark & \xmark\\
    \bottomrule
  \end{tabular}
\end{table}

\begin{itemize}
  \item \textbf{RotatE} row (slide~60): score $-\norm{\vh\circ\vr-\vt}$ with
    $\vh,\vt,\vr\in\C^k$; $\circ$ is the \textbf{Hadamard product}, e.g.\
    $\left[\begin{smallmatrix}3&5&7\\4&9&8\end{smallmatrix}\right]\circ
    \left[\begin{smallmatrix}1&6&3\\0&2&9\end{smallmatrix}\right]=
    \left[\begin{smallmatrix}3\times1&5\times6&7\times3\\
      4\times0&9\times2&8\times9\end{smallmatrix}\right]$.
  \item \textbf{TransE and RotatE:} they both satisfy a \textbf{weaker notion of
    1-to-N} -- that many tails can be equidistant to $\vr\ast\vh$ / $\vr+\vh$.
\end{itemize}

\subsection{KG embeddings in practice}\label{sec:lec10-practice}
\slides{61--62}

\begin{enumerate}
  \item Different KGs may have \textbf{drastically different relation
    patterns}!
  \item There is not a general embedding that works for all KGs, use the
    \textbf{table} (\cref{tab:lec10-patterns}) to select models.
  \item Try \textbf{TransE} for a quick run if the target KG does not have much
    symmetric relations.
  \item Then use more expressive models, e.g.\ \textbf{ComplEx}, \textbf{RotatE}
    (TransE in complex space; Sun et al., \emph{RotatE: Knowledge Graph Embedding
    by Relational Rotation in Complex Space}, ICLR 2019).
\end{enumerate}
\begin{itemize}
  \item \textbf{Empirical comparison} on FB15k-237 and WN18RR:
    \cref{tab:lec10-results} (MR = mean rank, MRR = mean reciprocal rank, H10 =
    Hits@10; $\downarrow$ lower / $\uparrow$ higher is better). RotatE (bold on
    the slide) is best on every metric; no TransR numbers on the slide.
\end{itemize}

\begin{table}[htbp]
  \centering\small
  \caption{Empirical comparison (slide~62).}\label{tab:lec10-results}
  \begin{tabular}{@{}lcccccc@{}}
    \toprule
    & \multicolumn{3}{c}{FB15k-237} & \multicolumn{3}{c}{WN18RR}\\
    \cmidrule(lr){2-4}\cmidrule(l){5-7}
    Model & MR$\downarrow$ & MRR$\uparrow$ & H10$\uparrow$
          & MR$\downarrow$ & MRR$\uparrow$ & H10$\uparrow$\\
    \midrule
    TransE   & 357 & .294 & .465 & 3384 & .226 & .501\\
    TransR   & --  & --   & --   & --   & --   & --\\
    DistMult & 254 & .241 & .419 & 5110 & .43  & .49\\
    ComplEx  & 339 & .247 & .428 & 5261 & .44  & .51\\
    RotatE   & \textbf{177} & \textbf{0.338} & \textbf{0.533}
             & \textbf{3340} & \textbf{0.476} & \textbf{0.571}\\
    \bottomrule
  \end{tabular}
\end{table}

\subsection{Summary}\label{sec:lec10-summary}
\slides{63}

\begin{itemize}
  \item \textbf{Link prediction / graph completion} is one of the prominent tasks
    on knowledge graphs.
  \item TransE / TransR / DistMult / ComplEx: models with different
    \textbf{embedding space} and \textbf{expressiveness}.
  \item \textbf{Next:} reasoning in knowledge graphs.
\end{itemize}

% ---------------------------------------------------------------------
\bigskip\phantomsection
\addcontentsline{toc}{section}{Lecture 10 at a glance}
\begin{summarybox}[Lecture 10 at a glance]
  \begin{itemize}
    \item \textbf{KG:} heterogeneous graph of entities (typed nodes) and
      relationships; massive and incomplete (Freebase: 93.8\% of persons without
      place of birth). \textbf{KG completion:} given (head, relation), predict
      missing tails.
    \item \textbf{Shallow embeddings} of entities and relations (no GNN); triple
      $(h,r,t)$: embedding of $(h,r)$ close to that of $t$; score $f_r(h,t)$ high
      iff $(h,r,t)$ exists.
    \item \textbf{Relation patterns:} symmetric / antisymmetric, inverse
      $r_2(h,t)\Rightarrow r_1(t,h)$, composition
      $r_1(x,y)\wedge r_2(y,z)\Rightarrow r_3(x,z)$, 1-to-N.
    \item \textbf{TransE:} $\vh+\vr\approx\vt$, $f_r=-\norm{\vh+\vr-\vt}$; margin
      loss $\sum[\gamma+d(\vh+\vr,\vt)-d(\vh'+\vr,\vt')]_+$ with corrupted
      triplets; no symmetric, no 1-to-N.
    \item \textbf{TransR:} $\vh_\perp=\mM_r\vh$, $\vt_\perp=\mM_r\vt$,
      $f_r=-\norm{\vh_\perp+\vr-\vt_\perp}$; all five patterns (composition via
      $\Ker(\mM_{r_3})=\Ker(\mM_{r_1})+\Ker(\mM_{r_2})$).
    \item \textbf{DistMult:} $f_r=\langle\vh,\vr,\vt\rangle=\sum_i\vh_i\vr_i\vt_i$
      (diagonal bilinear; cosine of $\vh\cdot\vr$ and $\vt$); symmetric and
      1-to-N only.
    \item \textbf{ComplEx:} $f_r=\cRe(\sum_i\vh_i\vr_i\bar\vt_i)$ in $\C^k$;
      adds antisymmetric and inverse ($\vr_1=\bar\vr_2$); no composition.
    \item \textbf{In practice:} no embedding works for all KGs; pick by the table;
      TransE for a quick run, then ComplEx / RotatE; RotatE best on FB15k-237 and
      WN18RR.
  \end{itemize}
\end{summarybox}
